\section{Reconstruction, mass gap, and non-Gaussian correlations}
\label{sec:os}

Section~\ref{sec:continuum} constructed the plane Schwinger functions
$S_n$ and proved their Euclidean covariance and growth bounds.  We now
verify the remaining reconstruction hypotheses and transfer the uniform
lattice decay estimate to the entire reconstructed Hilbert space.  A
separated four-point test will then show that the limiting field does not
satisfy Wick's rule.

Write $x=(x^0,x^1)$, with $x^0$ the Euclidean time coordinate, and put
$\theta(x^0,x^1)=(-x^0,x^1)$.  We use $\phi^a(f)$, $a\in\{1,2,3\}$, for
the limiting smeared fields supplied by
Proposition~\ref{prop:continuum-limit}.  Expectations of polynomials in
these variables are the corresponding pairings with the $S_n$.
For $s\in\R$, set
\[
 (\tau_s f)(x^0,x^1)=f(x^0-s,x^1),
 \qquad (\theta f)(x)=f(\theta x).
\]
The same notation $\tau_s$ denotes the induced translation of a field
polynomial.  Define $\Theta\phi^a(f)=\phi^a(\overline{\theta f})$ and
extend $\Theta$ multiplicatively and antilinearly to field polynomials.

\subsection{Reflection positivity and exponential clustering}

\begin{lemma}[Reflection positivity]
\label{lem:reflection-positivity}
Let $F$ be a polynomial in finitely many component fields
$\phi^a(f)$, where the smooth compactly supported tests $f$ have support
in $\{x^0>0\}$.  Then
\begin{equation}
 \E\bigl[(\Theta F)F\bigr]\ge0.
 \label{eq:continuum-reflection-positivity}
\end{equation}
This holds jointly for all three components.  The Schwinger functions
also satisfy the reality relations for a Hermitian multiplet and
$S_0=1$.
\end{lemma}

\begin{proof}
At every cutoff the nearest-neighbour measure is reflection positive
across microscopic site and link seams.  Since the union of the supports
appearing in $F$ is a compact subset of the open positive half-plane, a
microscopic seam within $O(a_N)$ of $\{x^0=0\}$ separates these supports
from their reflections for all sufficiently large $N$.  Evaluate $F$ on
the renormalized lattice fields and reflect about this seam.  Its
reflection-positive quadratic form is nonnegative.  The moment
convergence of Proposition~\ref{prop:continuum-limit} applies to each of
the finitely many terms in that form.  Replacing the microscopic
reflection by $\theta$ changes the tests by a translation of size
$O(a_N)$; Proposition~\ref{prop:moment-majorants} bounds the resulting
moment errors, which tend to zero.  Taking the limit proves
\eqref{eq:continuum-reflection-positivity}.  This argument permits
arbitrary component indices and polynomial coefficients, so it proves
joint positivity rather than separate componentwise positivity.
Reality and normalization pass through the same moment limits.
\end{proof}

The next estimate retains an exponent independent of the polynomial.
That uniformity will exclude low-energy states throughout the Hilbert
space, even though the constants multiplying the exponential depend on
the observables.

\begin{proposition}[Polynomial slab clustering]
\label{prop:slab-clustering}
There is $m>0$ with the following property.  Let $F$ and $G$ be fixed
polynomials in finitely many smooth compactly supported smeared
component fields.  If time translation separates the slabs containing
their supports by a distance $R\ge0$, then
\begin{equation}
 \bigl|\Cov(F,\tau_sG)\bigr|
       \le C_{F,G}e^{-mR}.
 \label{eq:polynomial-slab-clustering}
\end{equation}
The exponent $m$ is common to all such polynomials.  The Schwinger
functions satisfy the Euclidean clustering axiom also for Schwartz
tests.
\end{proposition}

\begin{proof}
Let $X_{i,N}$ denote the real smeared lattice fields entering the two
polynomials, with the translations needed to place their slabs.  Complex
tests may first be split into real and imaginary parts.  The factorial
moment bounds of Proposition~\ref{prop:moment-majorants}, uniformly in
$N$ and in translations of each fixed test, imply that for some
$\varepsilon>0$
\begin{equation}
 \sup_N \E\exp\Bigl(\varepsilon\sum_i|X_{i,N}|\Bigr)<\infty.
 \label{eq:polynomial-exponential-tails}
\end{equation}
Indeed, the even moment bounds give factorial bounds on all absolute
moments by Cauchy--Schwarz.  Summing the exponential series at a
sufficiently small radius gives an exponential moment for each variable;
H\"older's inequality combines the finitely many variables.  Neither
radius depends on the translation of a test.

For $M\ge1$, let $T_M(u)=\max(-M,\min(u,M))$ and form $F_{N,M}$ and
$G_{N,M}$ by replacing every real smeared field by its truncation
$T_M(X_{i,N})$.  If the polynomial degrees are $d_F,d_G$, then
\[
 \norm{F_{N,M}}_\infty\le C_F(1+M)^{d_F},
 \qquad
 \norm{G_{N,M}}_\infty\le C_G(1+M)^{d_G}.
\]
Equation~\eqref{eq:polynomial-exponential-tails} also gives
\begin{equation}
 \norm{F_N-F_{N,M}}_{L^2}
 +\norm{G_N-G_{N,M}}_{L^2}
 \le C_{F,G}e^{-c_{F,G}M},
 \label{eq:polynomial-truncation}
\end{equation}
with uniformly bounded $L^2$ norms for $F_N,G_N$.  To obtain this bound,
expand the polynomial difference into terms supported on
$\{\max_i|X_{i,N}|>M\}$.  Each term is bounded by a fixed polynomial in
$\sum_i|X_{i,N}|$; the exponential moment absorbs its square and the tail
indicator.  Cauchy--Schwarz then shows that replacing the covariance by
that of the truncated variables costs at most
$C_{F,G}e^{-c_{F,G}M}$.

The truncated variables remain bounded local lattice observables in the
same slabs.  Equation~\eqref{eq:orig-26} supplies a fixed physical decay
rate $\mu>0$.  The number of microscopic edges between the slabs,
multiplied by $a_N$, differs from $R$ by at most $O(a_N)$; this rounding
cost is bounded uniformly.  Consequently
\[
 |\Cov(F_N,\tau_sG_N)|
 \le C_{F,G}(1+M)^{d_F+d_G}e^{-\mu R}
       +C_{F,G}e^{-c_{F,G}M}.
\]
Set $M=1+R^2$.  The polynomial prefactor in the first term is absorbed
by $e^{\mu R/2}$, with a constant depending on the degrees.  The second
term is also at most a test-dependent constant times $e^{-\mu R/2}$.
Thus the estimate holds with the common choice $m=\mu/2$.  Moment
convergence passes it to the limiting fields.

Euclidean rotation invariance, proved in
Proposition~\ref{prop:euclidean-regularity}, gives clustering when one
compact collection of insertions is translated to infinity in any
direction.  The constants can be kept uniform as that direction varies:
the rotated tests have uniformly bounded suprema and supports in a
common bounded set, so the unit-box moment bounds give a common radius
in \eqref{eq:polynomial-exponential-tails}.  Their polynomial degrees
and coefficients are unchanged.  To extend this assertion to Schwartz tests, first approximate
them by compactly supported product tests.  At every fixed order, the
product-of-unit-box majorants used to prove
Equation~\eqref{eq:orig-35} bound the approximation errors uniformly in
the translation: place the polynomial weights separately around the
origins of the two translated collections.  Finite sums of product
tests are dense in the corresponding Schwartz test spaces.  The same
uniform bounds therefore allow approximation first and translation to
infinity second, proving the full clustering axiom.
\end{proof}

\subsection{Osterwalder--Schrader reconstruction and the full gap}

\begin{proposition}[Relativistic reconstruction]
\label{prop:os-reconstruction}
The Schwinger functions $S_n$ reconstruct a tempered Wightman theory in
$1+1$ dimensions with a three-component Hermitian scalar field.  Its
Hilbert space $\mathcal H$ is positive, its fields are local, and its
unitary representation of the Poincar\'e group satisfies the spectrum
condition.  The vacuum $\Omega$ is unique and normalized.  In the
positive-time reflection realization of $\mathcal H$, Euclidean time
translations induce a strongly continuous semigroup
$e^{-sH_{\rm phys}}$, $s\ge0$, with $H_{\rm phys}\ge0$ and
$H_{\rm phys}\Omega=0$.
\end{proposition}

\begin{proof}
We apply the Osterwalder--Schrader reconstruction theorem under the
linear growth condition~\cite[Section~IV.1]{OS75}.  We specify its regularity
hypothesis to make clear that the distributional construction in
Section~\ref{sec:continuum} suffices.  For $F$ on $(\R^2)^n$, use the
Schwartz seminorm
\[
 |F|_p=\max_{|\alpha|\le p}\sup_x
       (1+|x|^2)^{p/2}|\partial^\alpha F(x)|.
\]
The linear growth condition requires a bound
$|S_n(F)|\le \sigma_n|F|_{sn}$ with fixed $s$ and
$\sigma_n\le A(n!)^D$ for fixed $A,D$.  Equation~\eqref{eq:orig-35}
gives this with $s=4$: its weighted supremum is the zero-derivative part
of $|F|_{4n}$, and $C^n n!\le A(n!)^D$ after increasing $A,D$.

Restricting to Schwartz tests flat on the coincidence diagonals
(all derivatives vanish whenever $x_i=x_j$) gives precisely the test
class defined in \cite[Section~II, following Equation~(2.1), p.~284]{OS75}.
Normalization and reality follow
from Lemma~\ref{lem:reflection-positivity}; permutation symmetry and
Euclidean invariance follow from
Proposition~\ref{prop:euclidean-regularity}; joint positive-time
reflection positivity follows from the same lemma, extended by
test-space continuity; and clustering is
Proposition~\ref{prop:slab-clustering}.  These are the normalized
Euclidean axioms with the required linear growth bound.

The theorem applies to a finite Hermitian multiplet with
scalar spacetime transformation law; see \cite[Section~I, Remark~5]{OS75}.
More explicitly, the test
functions carry indices $a_1,\ldots,a_n\in\{1,2,3\}$; reflection
positivity holds for arbitrary combinations of these indices, while
summing the component bounds costs at most an exponential factor in
$n$, already allowed by the preceding factorial estimate.  Thus the
multilinear reconstruction has a single positive Hilbert space and the
asserted three-component field.  Clustering gives uniqueness of the
vacuum.  The positive-time reflection realization and its translation
semigroup are part of this reconstruction.  This application uses
smeared distributions throughout and requires neither Euclidean point
fields nor a prescription for equal-time products.
\end{proof}

The uniform exponent in Proposition~\ref{prop:slab-clustering} now
controls every state generated by time-ordered fields.  Positivity of
spectral measures extends the resulting exclusion of low energies to
the Hilbert-space closure.  This is the same spectral implication as
in \Rref{prop:continuum}; we give the argument for the present
reconstructed states.

\begin{proposition}[Gap above the vacuum]
\label{prop:continuum-gap}
For the Hamiltonian in Proposition~\ref{prop:os-reconstruction}, the
vacuum complement obeys
\begin{equation}
 H_{\rm phys}\big|_{\Omega^\perp}\ge m>0,
 \label{eq:orig-36}
\end{equation}
where $m$ is the common exponent of
Proposition~\ref{prop:slab-clustering}.
\end{proposition}

\begin{proof}
We first justify the density of states to which the slab estimate
applies.  Coincidence-flatness does not imply flatness on a surface
where two times agree but the spatial coordinates differ.  The positive
position measures and Euclidean rotations nevertheless allow us to
remove these surfaces in the reflection norm.

Spacetime rotations act only on the insertion coordinates, since
the three internal components are spacetime scalars.  Thus each
fixed-index component measure of $S_k$ is separately rotation
invariant.  For any such measure $\mu_k$ and $i\ne j$, set
\[
 A_{ij}=\{x:x_i^0=x_j^0,\ x_i\ne x_j\}.
\]
Let $B$ be a ball centred at the origin in $(\R^2)^k$.  It has finite
$\mu_k$-mass by Proposition~\ref{prop:moment-majorants} and is
invariant under simultaneous spacetime rotations.  For each
configuration with $x_i\ne x_j$, only finitely many rotation angles
make the time component of $x_i-x_j$ vanish.  Rotation invariance and
Tonelli's theorem therefore give
\[
 \mu_k(A_{ij}\cap B)
  =\frac1{2\pi}\int_B\int_0^{2\pi}
       \mathbf 1_{A_{ij}}(O_\alpha x)\,\dd\alpha\,\dd\mu_k(x)=0.
\]
Here $O_\alpha$ acts on every insertion coordinate.  Increasing the
ball shows that each equal-time surface has measure only on its full
coincidence diagonal.

The OS construction uses ordered positive-time tests; see
\cite[Section~V, following Equation~(5.2)]{OS75}.  One may also begin
with all coincidence-flat tests supported in positive time.  We show
that their reflection completion is the same.  Truncation at infinity
and away from the time-zero boundary reduces this comparison to
compactly supported tests in the latter class;
the smooth positive-time support condition and rapid decrease justify
these cutoffs in Schwartz topology, and the growth bounds imply
continuity of the reflection form.  Let $f_n$ be one of these compact
tests, with the coincidence-flatness used in reconstruction, and write
$[f_n]$ for its class for the reflection form.  Choose
a smooth even function $\chi_\varepsilon$ that is zero on
$[-\varepsilon,\varepsilon]$ and one outside
$[-2\varepsilon,2\varepsilon]$, with values in $[0,1]$, and put
\[
 f_{n,\varepsilon}(x)
   =f_n(x)\prod_{i<j}\chi_\varepsilon(x_i^0-x_j^0).
\]
The difference $f_n-f_{n,\varepsilon}$ tends to zero off the internal
equal-time surfaces.  Those surfaces have zero componentwise
$S_{2n}$-mass away from full coincidences, as just proved, and $f_n$
vanishes on the latter.  In the reflection form the positive and
reflected negative time arguments are separated from one another.
The integrands defining the squared reflection norm of this difference
thus tend to zero almost everywhere for every component measure of
$S_{2n}$.  They are bounded by a fixed constant on a common compact
set.  Dominated convergence proves
\[
 \norm{[f_n]-[f_{n,\varepsilon}]}\longrightarrow0.
\]
This approximation is in the reflection norm; it does not assert
Schwartz density after cutting equal-time surfaces.

The support of $f_{n,\varepsilon}$ lies in finitely many strict time
chambers.  Split it according to the ordering and use permutation
symmetry, permuting component indices with insertion coordinates, to
put each term in $0<x_1^0<\cdots<x_n^0$.  A compact subset
of this chamber is covered by finitely many product boxes whose time
intervals are disjoint and ordered.  A partition of unity and
product-test approximation in these boxes give finite sums of products
of smooth compactly supported one-field tests, converging in Schwartz
topology and hence in reflection norm.  Finite component sums obey the
same argument.  These compact time-ordered product states are dense in
the OS space $\mathcal H$, and the approximation also identifies the
completion of the larger coincidence-flat positive-time test class isometrically with
$\mathcal H$.

Let $F$ be a finite sum of products of smeared fields, with each product
supported in strictly ordered, mutually disjoint positive-time
intervals, as constructed above.  Its test kernel is coincidence-flat.
Write $[F]$ for its reflection-space class in $\mathcal H$ and set
\[
 v=[F]-\E[F]\Omega.
\]
These centered vectors are dense in $\Omega^\perp$, and their semigroup
matrix elements are
\[
 \inner{v}{e^{-sH_{\rm phys}}v}
   =\E\bigl[(\Theta F)(\tau_sF)\bigr]
       -\overline{\E[F]}\E[F].
\]
For all sufficiently large $s$, the reflected and translated supports
lie in slabs separated by $s$ plus a fixed constant depending on $F$.
Apply Proposition~\ref{prop:slab-clustering} with first observable
$\overline{\Theta F}$ and second observable $F$, using the
sesquilinear covariance convention of Section~\ref{sec:trace}.  It implies
\begin{equation}
 0\le\inner{v}{e^{-sH_{\rm phys}}v}\le C_v e^{-ms}.
 \label{eq:os-vector-decay}
\end{equation}
For a finite sum of product states one applies the same estimate to
each of the finitely many cross terms; the exponent remains $m$.

Let $E_H$ be the spectral resolution of $H_{\rm phys}$ and
$\nu_v(B)=\inner{v}{E_H(B)v}$ its finite positive spectral measure.
For $0<m'<m$, positive measure in $[0,m')$ would give
\[
 \inner{v}{e^{-sH_{\rm phys}}v}
    =\int_{[0,\infty)}e^{-s\lambda}\,\dd\nu_v(\lambda)
    \ge e^{-sm'}\nu_v([0,m')),
\]
contradicting \eqref{eq:os-vector-decay} as $s\to\infty$.
Hence $E_H([0,m'))v=0$ for every centered vector in the dense set.
Since a spectral projection is bounded, it vanishes on all of
$\Omega^\perp$.  Taking $m'\uparrow m$ proves
\eqref{eq:orig-36}.
\end{proof}

\subsection{Surviving fields and a separated fourth cumulant}

It remains to verify that the vacuum complement is nonzero and that
the limiting field is non-Gaussian.  Both follow from observables at
the fixed terminal block scale.  This uses the same field
renormalization as the construction and the spectral estimate.

\begin{proposition}[Nontriviality and failure of Wick factorization]
\label{prop:non-gaussian}
For a sufficiently large fixed terminal coupling $H$, there exist four
smooth compactly supported real tests $f_1,\ldots,f_4$ with strictly
disjoint ordered time ranges for which
\begin{align}
 &\E\prod_{i=1}^4\phi^1(f_i)
 -\E[\phi^1(f_1)\phi^1(f_2)]\E[\phi^1(f_3)\phi^1(f_4)] \notag\\
 &\quad
 -\E[\phi^1(f_1)\phi^1(f_3)]\E[\phi^1(f_2)\phi^1(f_4)]
 -\E[\phi^1(f_1)\phi^1(f_4)]\E[\phi^1(f_2)\phi^1(f_3)]
 \ne0.
 \label{eq:nonzero-fourth-cumulant}
\end{align}
The reconstructed Hilbert space satisfies $\Omega^\perp\ne\{0\}$.
Consequently the spectral bottom on $\Omega^\perp$ is both positive
and finite.
\end{proposition}

\begin{proof}
Choose four distinct unit terminal cells at fixed bounded distances,
with gaps between their closed time intervals.  Their number,
placement, and separations are fixed independently of $H$ and of the
cutoff.  Translations of the block partitions, as allowed in
Section~\ref{sec:continuum}, realize this common physical placement.
Let $X_{i,N}$ be the first component of the renormalized field sum over
the descendants of cell $i$.

The local source estimate of Proposition~\ref{prop:source-step} and
the comparison in Equation~\eqref{eq:orig-34} give, on a fixed complex
neighbourhood of the origin,
\begin{equation}
 \E\exp\Bigl(\sum_{i=1}^4 z_iX_{i,N}\Bigr)
   =\E_\mu\exp\Bigl(\sum_{i=1}^4 z_iV_{Y_i}^1\Bigr)+o_H(1),
 \label{eq:cell-generating-comparison}
\end{equation}
uniformly in the cutoff and in the admitted axis tori.  Here $\mu$ is
the corresponding terminal spin law and $Y_i$ labels cell $i$.
Cauchy's formula transfers the uniform analytic error to each of the
fixed moments of order at most four.

For completeness, the alignment estimate
\eqref{eq:orig-23} controls the spins at these finitely many sites in
mean square.  If a nearest-neighbour path of length $\ell$ joins $Y_i$
to $Y_1$, then
\[
 \E_\mu|V_{Y_i}-V_{Y_1}|^2
 \le \ell\sum_{e\text{ on the path}}\E_\mu r_e^2
 \le \ell^2\bigl(Ct^2+4e^{-cHt^2}\bigr)=o_H(1).
\]
The path lengths are fixed; $t=t_0$ tends to zero and $Ht^2=p_0^2$
tends to infinity as $H\to\infty$.  Internal $O(3)$ invariance makes
each $V_Y$ uniformly distributed on $S^2$.  If $U$ has this law, then
\[
 \E U^1=0,\qquad \E(U^1)^2=\frac13,
 \qquad \E(U^1)^4=\frac15.
\]
The coordinates have absolute value at most one.  Telescoping a
product and using the preceding mean-square estimate therefore gives,
for every pair of distinct chosen cells and for the four-cell product,
\begin{equation}
 \E[X_{i,N}X_{j,N}]=\frac13+o_H(1),
 \qquad
 \E\prod_{i=1}^4X_{i,N}=\frac15+o_H(1).
 \label{eq:terminal-cell-moments}
\end{equation}
The first moments vanish exactly by internal symmetry.

The estimates are uniform on axis tori.  We may thus first take the
periodic infinite-plane limit at fixed cutoff and then the continuum
limit.  The descendants in a standard blocking are axis boxes;
Lemma~\ref{lem:boundary-strips} identifies their limits with the fixed
physical cells and justifies any converging displacement of their
boundaries.  Denote the resulting cell variables by $X_i$.  Their
fourth cumulant is
\begin{equation}
 \E\prod_{i=1}^4X_i
   -\sum_{\{\{i,j\},\{k,l\}\}}\E[X_iX_j]\E[X_kX_l]
   =\frac15-3\left(\frac13\right)^2+o_H(1)
   =-\frac{2}{15}+o_H(1),
 \label{eq:cell-fourth-cumulant}
\end{equation}
where the sum runs over the three pairings of $\{1,2,3,4\}$.
Choose $H$ sufficiently large that this quantity is nonzero.  This
choice depends only on the fixed cells and the earlier RG constants,
and can be made before the later choices of $M_0$ and $K_0$ in
Section~\ref{sec:trace}.

Approximate each cell indicator by bounded smooth tests whose
differences from the indicator are supported in shrinking boundary
strips.  The approximations can preserve the strict gaps between the
time ranges.  Lemma~\ref{lem:boundary-strips} makes every moment in
\eqref{eq:cell-fourth-cumulant} converge under this approximation.
Thus suitable smooth tests satisfy
\eqref{eq:nonzero-fourth-cumulant}.  Since they have separated ordered
times, these are Schwinger correlations in the domains used by
reconstruction.  Wick factorization of the reconstructed Wightman
functions would, by their Euclidean continuation, give Wick
factorization for these Schwinger correlations as well.  Equation
\eqref{eq:nonzero-fourth-cumulant} excludes it.

The same two-cell calculation can be made for unit cells reflected
across $x^0=0$, with both cells a positive distance from the seam.
Their first-component correlation is $1/3+o_H(1)>0$.  Apply reflected
smooth approximations to this pair.  For a suitable real test $f$
supported strictly in positive time,
\[
 \norm{[\phi^1(f)]}^2
   =\E\bigl[\phi^1(\theta f)\phi^1(f)\bigr]>0,
 \qquad
 \inner{\Omega}{[\phi^1(f)]}=\E\phi^1(f)=0.
\]
Hence $\Omega^\perp$ contains a nonzero vector $v$.  Its spectral
measure has total mass $\norm v^2>0$.  Since $[0,\infty)$ is the
increasing union of bounded intervals, some finite $E$ satisfies
$\nu_v([0,E])>0$.  The spectral bottom on $\Omega^\perp$ is therefore
finite, and Proposition~\ref{prop:continuum-gap} makes it positive.
\end{proof}

\begin{proof}[Completion of the proof of Theorem~\ref{thm:main}]
The shooting construction gives bare couplings $\beta_N\to\infty$.
The prescribed lattice spacings satisfy $a_N=L^{-N}\to0$, and the
source construction in Section~\ref{sec:sources} supplies the positive
field renormalizations $B_N$ for the standard nearest-neighbour $O(3)$
model.  The terminal kinetic coupling is the fixed number $H$ at unit
physical block length.
Proposition~\ref{prop:continuum-limit} establishes convergence of the
renormalized correlations and finite lists of smeared fields, with
infinite volume taken first or simultaneously along the specified torus
exhaustion.  Proposition~\ref{prop:euclidean-regularity} gives their
temperedness and Euclidean symmetries.

Proposition~\ref{prop:os-reconstruction} reconstructs the local,
unitary relativistic theory.  Its full Hamiltonian gap is
Proposition~\ref{prop:continuum-gap}, and
Proposition~\ref{prop:non-gaussian} proves that a nonzero vacuum
complement and a nonvanishing connected fourth correlation survive in
this same continuum limit.  In particular the physical gap is positive
and finite at the fixed scale of the construction.  Every limiting
observable uses the field renormalizations already prescribed from the
nearest-neighbour model; the blocking observations are integrated
exactly and introduce no additional microscopic physical fields.  The
construction adds no symmetry-breaking mass term.
\end{proof}
